\begin{abstract}
A fast learned codec still applies the same synthesis depth to every image
region. FLEX-UF studies a different allocation of computation: spend depth
where it improves reconstruction, and stop elsewhere. We develop this
principle in the intra model of DCVC-UF through spatial early exit. One
coded representation feeds a shared stem and a nested synthesis path;
tiles depart through exit adapters before a common repair and reconstruction
head. This mechanism connects naturally to a complementary, ClassSR-inspired
design that selects among independently trained 2/4/6/8/10/12-block codecs.
We specify that controlled depth study while reporting completed results
only for the shared-exit path. On 53 intra frames at five quality indices,
learned routing removes \MainRouterSaving\% of modelled synthesis
multiply--accumulates at a nominal 0.1\,dB target, including \MainPremium\
percentage points beyond ordered dithering. Retrospective selection under a
common delivered RGB-loss cap retains a 2.93-point margin over all 265 pairs.
The margin contracts as the cheapest exit saturates: adaptation is most
valuable when regions benefit differently from additional depth.
These results establish an arithmetic--quality trade-off under
source-calibrated control; autonomous calibration, model-bank routing and
complete-bitstream latency remain to be evaluated.
\end{abstract}
